\section*{अध्याय \ref{ch:b3:residues} --- लोरां श्रेणी और
अवशेष प्रमेय}

\begin{solution}{exo:b3:residues:1}
$\frac{\sin z}z = 1 - \frac{z^2}6 + \cdots$: विलोपनीय, \omterm{def:b3:residues:singularities}{अवशेष} $0$। $\frac{\eu^z - 1}{z^2} = \frac1z + \frac12 +
\frac z6 + \cdots$: सरल ध्रुव, \omterm{def:b3:residues:singularities}{अवशेष} $1$। $\frac1{z(z-1)^2}$: $0$ पर
\omterm{def:b3:residues:singularities}{अवशेष} $\frac1{(0-1)^2} = 1$ वाला सरल ध्रुव; और $1$ पर \omterm{def:b3:residues:singularities}{अवशेष} $\frac{\dd}{\dd z}\bigl(\frac1z\bigr)\big|_{z=1} = -1$ वाला द्विक ध्रुव।
$\frac{\cos z}{z^3} = \frac1{z^3} - \frac1{2z} + \cdots$: कोटि $3$ का ध्रुव, \omterm{def:b3:residues:singularities}{अवशेष} $-\frac12$। $\eu^{1/z} = \sum_{n\geq0}
\frac{z^{-n}}{n!}$: सारभूत, \omterm{def:b3:residues:singularities}{अवशेष} $1$।
$\frac1{\sin z}$: $n\pi$ पर सरल ध्रुव; $0$ पर \omterm{def:b3:residues:singularities}{अवशेष} $\frac1{\cos 0} = 1$, और $\pi$ पर $\frac1{\cos\pi} = -1$ (\cref{met:b3:residues:compute},
$g/h'$)।
\end{solution}

\begin{solution}{exo:b3:residues:2}
$z = \eu^{\iu t}$, $\cos t = \frac{z + z^{-1}}2$, $\dd t
= \frac{\dd z}{\iu z}$ के साथ:
\[
\int_0^{2\pi}\frac{\dd t}{a + \cos t}
= \oint_{\abs z = 1}\frac{2\,\dd z}{\iu\,(z^2 + 2az + 1)} .
\]
मूल $z_\pm = -a \pm \sqrt{a^2 - 1}$ $\abs{z_+} < 1 < \abs{z_-}$ के साथ $z_+z_- =
1$ संतुष्ट करते हैं; $z_+$ पर \omterm{def:b3:residues:singularities}{अवशेष} $\frac{1}{z_+ - z_-} = \frac1{2\sqrt{a^2-1}}$ है,
अतः समाकल $\frac2\iu\cdot2\iu\pi\cdot\frac1{2\sqrt{a^2-1}} =
\frac{2\pi}{\sqrt{a^2-1}}$ है। $a \to 1^+$ पर वह फट जाता है (समाकल्य $t = \pi$ पर शिखर बनाता
है); और $a \to \infty$ पर वह $\frac{2\pi}a$ की भाँति व्यवहार करता है, जो $\int\frac{\dd t}a$ से मेल
खाता है।
\end{solution}

\begin{solution}{exo:b3:residues:3}
पहला समाकल: $\frac{z^2}{1+z^6}$ के ऊपरी ध्रुव $p =
\eu^{\iu\pi/6}, \iu, \eu^{5\iu\pi/6}$ पर; और प्रत्येक पर $\operatorname{Res} = \frac{p^2}{6p^5} = \frac{p^3}{6p^6} =
-\frac{p^3}6$, तथा
$p^3$ मान $\iu, -\iu, \iu$ लेता है: अतः अवशेषों का योग $-\frac{\iu}{6}$। चाप $O(R^{-4})\cdot O(R) \to 0$ है:
\[
\int_\R\frac{x^2\,\dd x}{1 + x^6} =
2\iu\pi\Bigl(-\frac \iu6\Bigr) = \frac\pi3 .
\]
दूसरा: $\frac1{(1+z^2)^2} =
\frac1{(z-\iu)^2(z+\iu)^2}$ का $\iu$ पर द्विक ध्रुव:
\[
\operatorname{Res} = \frac{\dd}{\dd z}\,(z +
\iu)^{-2}\Big|_{z=\iu} = \frac{-2}{(2\iu)^3} =
\frac{-2}{-8\iu} = -\frac\iu4,
\qquad
\int_\R\frac{\dd x}{(1+x^2)^2} =
2\iu\pi\cdot\Bigl(-\frac\iu4\Bigr) = \frac\pi2 ,
\]
जो \cref{exo:b3:fouriertransform:3}(ख) के संगत है।
\end{solution}

\begin{solution}{exo:b3:residues:4}
$\frac{\eu^{\iu tz}}{z^2 + a^2}$ को ऊपरी अर्ध-समतल में बंद कीजिए ($t \geq 0$): वहाँ $\abs{\eu^{\iu tz}} =
\eu^{-t\operatorname{Im}z} \leq 1$, अतः चाप का
योगदान $O(R^{-2})\cdot O(R) \to 0$ है। एकमात्र घिरा हुआ ध्रुव $\iu a$ \omterm{def:b3:residues:singularities}{अवशेष} $\frac{\eu^{-at}}{2\iu a}$ के साथ सरल
है:
\[
\int_\R\frac{\eu^{\iu tx}}{x^2 + a^2}\dd x =
\frac{\pi}{a}\,\eu^{-at},
\qquad\text{अतः}\qquad
\int_\R\frac{\cos(tx)}{x^2+a^2}\dd x = \frac\pi
a\,\eu^{-a\abs t}
\]
(वास्तविक भाग; $t$ में सम)। यह $\widehat{\eu^{-a\abs x}} = \frac{2a}{a^2+\xi^2}$ (\cref{exo:b3:fouriertransform:1}) का प्रतिलोमन प्रतिरूप
है: दोनों परिकलन \cref{thm:b3:fouriertransform:inversion} के माध्यम से एक-दूसरे की पुष्टि करते हैं।
\end{solution}

\begin{solution}{exo:b3:residues:5}
आंशिक भिन्न: $f = \frac1{z-2} - \frac1{z-1}$। $\abs z < 1$ पर (टेलर): $f = \sum_{n\geq0}\bigl(1 - 2^{-n-1}\bigr)z^n$। $1 < \abs z < 2$ पर: $\frac1{z-2} =
-\sum_{n\geq0}\frac{z^n}{2^{n+1}}$ और $-\frac1{z-1} =
-\sum_{n\geq1}z^{-n}$:
अर्थात् एक सच्ची दो-मुखी श्रेणी। $\abs z > 2$ पर: $f = \sum_{n\geq1}\bigl(2^{n-1} -
1\bigr)z^{-n}$। तीनों भिन्न हैं
क्योंकि लोरां प्रसार \emph{वलयों} से जुड़े होते हैं, बिंदुओं से नहीं:
प्रत्येक क्षेत्र के अपने ज्यामितीय प्रसार होते हैं। $z^{-1}$-गुणांक
(क्रमशः $-1$ और $0$) \emph{भीतर की एकलताओं} को घेरने वाले वृत्तों
पर समाकल हैं, न कि $0$ पर \omterm{def:b3:residues:singularities}{अवशेष} ($f$ $0$ पर \omterm{def:b3:holomorphic:holo}{समविश्लेषिक} है):
$1 < \abs z < 2$ के लिए गुणांक $-1$ $\operatorname{Res}(f, 1)$ है; और $\abs z > 2$ के लिए गुणांक $0$ $\operatorname{Res}(f,1) +
\operatorname{Res}(f,2) = -1 + 1$ है।
$f$ के \omterm{def:b3:residues:singularities}{अवशेष}: $1$ पर $-1$ और $2$ पर $+1$।
\end{solution}

\begin{solution}{exo:b3:residues:6}
(क) \omterm{thm:b3:residues:laurent}{लोरां श्रेणी} $\sum_nz^{-n}/n!$ में अनंत ऋणात्मक पद हैं: अतः सारभूत। $\eu^{1/z} = w$
($w \neq
0$) हल करने पर: $\frac1z = \log\abs w + \iu\arg w + 2\iu\pi k$, अतः
\[
z_k = \frac1{\log\abs w + \iu\arg w + 2\iu\pi k}
\xrightarrow[k\to\infty]{} 0 :
\]
अर्थात् प्रत्येक शून्येतर मान $0$ के निकट अनंत बार प्राप्त होता है
--- जो सघनता से प्रबलतर है।
(ख) $z = 1/x$ के अनुदिश $x \to +\infty$: $\eu^x \to \infty$; और $z = \iu/y$ के अनुदिश: मापांक $1$।
किसी ध्रुव के लिए \emph{प्रत्येक पहुँच के अनुदिश} $\abs{f(z)} \to \infty$ चाहिए: यहाँ
सीमा $[0, +\infty]$ में भी विद्यमान ही नहीं है।
\end{solution}

\begin{solution}{exo:b3:residues:7}
(क) $\abs z = 1$ पर: $\abs{z^7 + z - 1} \leq 3 < 4 =
\abs{-4z^3}$। $f = -4z^3$, $g = z^7 + z - 1$ के साथ रूशे: अतः चक्र में तीन
शून्य।
(ख) $\abs z = 1$ पर: $\abs{z^4 + 1} \leq 2 < 5 = \abs{5z}$: अतः $\abs z < 1$ में एक शून्य। $\abs z = 2$ पर: $\abs{5z + 1} \leq
11 < 16 = \abs{z^4}$: अतः
$\abs z < 2$ में चार शून्य। इसलिए वलय में तीन शून्य।
(ग) $P = z^n + a_{n-1}z^{n-1} + \dots$ के लिए: $\abs z = R >
1 + \sum\abs{a_k}$ पर $\abs{P - z^n} \leq
\bigl(\sum\abs{a_k}\bigr)R^{n-1} < R^n = \abs{z^n}$: अतः $P$ के $D(0, R)$ में ठीक $n$
शून्य हैं --- अर्थात् बहुलता सहित दालांबेर--गाउस, विशुद्ध गणना से।
\end{solution}

\begin{solution}{exo:b3:residues:8}
(क) मान लीजिए $f(a) = 0$, $f \not\equiv 0$: तो ऐसा $r$ चुनिए कि $f$ वृत्त $C = \partial D(a, r)$
पर शून्य-रहित हो (\omterm{thm:b3:holomorphic:identity}{विविक्त शून्य}) और $m = \min_C\abs f > 0$। \cref{thm:b3:holomorphic:weierstrassconv} से $C$ पर एकसमान
रूप से $f_n \to f$ और $f_n' \to f'$; और बड़े $n$ के लिए $C$ पर $\abs{f_n}
\geq m/2$, अतः
\[
\frac1{2\iu\pi}\int_C\frac{f_n'}{f_n}
\longrightarrow \frac1{2\iu\pi}\int_C\frac{f'}{f} \geq 1
\]
(सीमा शून्य $a$ को गिनती है; और अभिसरण इसलिए कि अंश एकसमान रूप से
अभिसरित होते हैं और हर एकसमान रूप से नीचे से परिबद्ध हैं)। बायाँ
पक्ष चक्र में $f_n$ के शून्य गिनने वाला कोई पूर्णांक है: अतः उसे
अंततः $\geq 1$ होना पड़ेगा --- जो शून्य-रहितता का खंडन करता है। इसलिए
$f$ शून्य-रहित है।

(ख) $w \in \Omega$ स्थिर कीजिए और \omterm{def:b3:topology:connected}{संबद्ध} \omterm{def:b3:topology:topology}{विवृत समुच्चय} $\Omega\setminus\{w\}$ ($\C$ के किसी
विवृत \omterm{def:b3:topology:connected}{संबद्ध} उपसमुच्चय से एक बिंदु हटाने पर \omterm{def:b3:topology:connected}{संबद्धता} बनी रहती है)
पर (क) को $g_n(z) = f_n(z) -
f_n(w)$ पर लगाइए, जो एकैकीपन से वहाँ शून्य-रहित है और $g = f
- f(w)$
पर अभिसरित होता है। यदि $f$ अचरेतर हो, तो $\Omega\setminus\{w\}$ पर $g \not\equiv 0$, अतः
$g$ वहाँ शून्य-रहित है: सभी $z \neq w$ के लिए $f(z) \neq
f(w)$। और $w$ स्वेच्छ
था, अतः $f$ एकैकी है।
\end{solution}

\begin{solution}{exo:b3:residues:9}
त्रिज्यखंड की परिसीमा $[0, R]$, चाप $A_R$, और उलटी गई किरण $\eu^{2\iu\pi/n}[0, R]$ से
बनी है। भीतर $\frac1{1+z^n}$ का एकमात्र ध्रुव $p = \eu^{\iu\pi/n}$ है, जिसका \omterm{def:b3:residues:singularities}{अवशेष} $\frac1{np^{n-1}} = \frac{p}{np^n} = -\frac pn$ है।
लौटती किरण पर $z = \eu^{2\iu\pi/n}x$ $z^n = x^n$ और $\dd z = \eu^{2\iu\pi/n}\dd x$ दे देता है; और चाप $O(R^{-n})\cdot O(R) \to 0$ है। अतः
\[
\bigl(1 - \eu^{2\iu\pi/n}\bigr)
\int_0^\infty\frac{\dd x}{1 + x^n}
= 2\iu\pi\Bigl(-\frac{\eu^{\iu\pi/n}}n\Bigr),
\quad\text{अतः}\quad
\int_0^\infty\frac{\dd x}{1+x^n}
= \frac{2\iu\pi}{n\,\bigl(\eu^{\iu\pi/n} -
\eu^{-\iu\pi/n}\bigr)} = \frac{\pi}{n\sin(\pi/n)} .
\]
$n = 2$: $\frac\pi{2\sin(\pi/2)} = \frac\pi2 =
[\arctan]_0^\infty$। $n \to \infty$ पर: मान $1$ की ओर जाता है, और वस्तुतः समाकल्य
$\mathbf 1_{\intco01}$ की ओर जाता है ($n \geq 2$ के लिए प्रभावी अभिसरण प्रमेय के प्रभुत्व
$\min(1, x^{-2})$ के साथ)।
\end{solution}

\begin{solution}{exo:b3:residues:10}
बड़े $\abs z = R$ पर $\abs f \leq C R^{-2}$: $\abs{\oint_{C_R}
f} \leq 2\pi R\cdot CR^{-2} \to 0$। पर सभी ध्रुवों से परे $R$ के लिए \omterm{def:b3:residues:singularities}{अवशेष}
प्रमेय $\oint_{C_R}f =
2\iu\pi\sum_{\text{सभी }p}\operatorname{Res}(f, p)$ देती है: अतः कुल योग लुप्त हो जाता है। $f = \frac1{z(z-1)(z-2)}$ के लिए:
$0$ पर \omterm{def:b3:residues:singularities}{अवशेष} $\frac1{(-1)(-2)} = \frac12$, $1$ पर $\frac1{1\cdot(-1)} = -1$, और $2$ पर $\frac1{2\cdot1} = \frac12$ --- जिनका
योग भविष्यवाणी के अनुसार $0$ है, और
\[
\frac1{z(z-1)(z-2)} = \frac{1/2}{z} - \frac1{z - 1} +
\frac{1/2}{z-2} :
\]
अर्थात् \omterm{def:b3:residues:singularities}{अवशेष} \emph{ही} आंशिक भिन्न गुणांक हैं, और शून्य-योग सर्वसमिका
मुफ़्त में संगति जाँच दे देती है (अथवा शेष से अंतिम गुणांक निर्धारित
कर देती है)।
\end{solution}

\begin{solution}{exo:b3:residues:11}
कुंजीछिद्र पर, चुने गए निर्धारण के साथ: कट के ठीक ऊपर $\log z = \ln x$; और ठीक
नीचे $\log z =
\ln x + 2\iu\pi$। चारों टुकड़े
\[
\Bigl(1 - \eu^{2\iu\pi(a-1)}\Bigr)\int_\varepsilon^R
\frac{x^{a-1}}{1+x}\dd x + \int_{C_R} + \int_{C_\varepsilon}
= 2\iu\pi\operatorname{Res}(f, -1) .
\]
देते हैं। चाप: $C_R$ पर $\abs{f} \leq \frac{R^{a-1}}{R - 1}$, लंबाई $2\pi R$: योगदान $O(R^{a-1}) \to 0$ ($a < 1$);
$C_\varepsilon$ पर $\abs f \leq \frac{\varepsilon^{a-1}}{1 - \varepsilon}$, लंबाई $2\pi\varepsilon$: $O(\varepsilon^a)
\to 0$ ($a > 0$)। \omterm{def:b3:residues:singularities}{अवशेष}: $z = -1 = \eu^{\iu\pi}$ पर $\operatorname{Res} = \eu^{(a-1)\iu\pi}$।
अतः
\[
I(a) = \frac{2\iu\pi\,\eu^{\iu\pi(a-1)}}{1 -
\eu^{2\iu\pi(a-1)}}
= \frac{2\iu\pi}{\eu^{-\iu\pi(a-1)} - \eu^{\iu\pi(a-1)}}
= \frac{\pi}{-\sin(\pi(a-1))} = \frac{\pi}{\sin\pi a} .
\]
गामा परावर्तन: $B(a, 1-a) =
\int_0^1t^{a-1}(1-t)^{-a}\dd t$; प्रतिस्थापन $t =
\frac x{1+x}$, $1 - t = \frac1{1+x}$, $\dd t =
\frac{\dd x}{(1+x)^2}$ उसे $\int_0^\infty
\frac{x^{a-1}}{1+x}\dd x = I(a)$ में
बदल देता है, और ऑयलर का सूत्र
$B(a, 1-a) = \Gamma(a)\Gamma(1-a)/\Gamma(1)$ 
(\cref{pb:b3:lebesgue:1}) $\Gamma(a)\Gamma(1-a) = \frac\pi{\sin\pi a}$ दे देता है --- विशेष रूप से एक बार फिर $\Gamma(\tfrac12) = \sqrt\pi$।

\end{solution}

\begin{solution}{exo:b3:residues:12}
(क) $\abs z = 1$ पर: $\abs{z^4} = 1 < 7 \leq \abs{8z + 1}$ ($\abs{8z} - 1 = 7$): अतः $P$ के $\mathbb D$ में उतने ही शून्य
हैं जितने $8z + 1$ के, अर्थात् \emph{एक} ($-\frac18$ पर)।

(ख) $\abs z = 3$ पर: $\abs{8z + 1} \leq 25 < 81 =
\abs{z^4}$: अतः $z^4$ के विरुद्ध रूशे चारों शून्य $\abs z < 3$ में
दे देते हैं, इसलिए वलय $1 <
\abs z < 3$ में $4 - 1 = 3$ शून्य। तीक्ष्णीकरण: $\abs z = r \geq 2.1$
वाला कोई शून्य $r^4 = \abs{8z + 1} \leq 8r + 1$ संतुष्ट करता, पर $r \geq 2$ के लिए $r^4 -
8r - 1$ वर्धमान है
और $r = 2.1$ पर $19.45 -
16.8 - 1 = 1.65 > 0$ के बराबर: असंभव। अतः तीनों बाहरी शून्य $1 < \abs z < 2.1$ में
हैं। (संख्यात्मक रूप से: $-1.95$ के निकट कोई वास्तविक शून्य और $1.04 \pm
1.73\iu$
के निकट कोई संयुग्म युग्म, जिसका मापांक $2.02$ है --- और इसीलिए ठीक
त्रिज्या $2$ पर रूशे का प्रयास विफल होना ही चाहिए: प्रमेय कड़ा
प्रभुत्व माँगती है, और शून्य ठीक बाहर बैठे हैं।)

(ग) $\iu\R$ पर कोई शून्य नहीं: $\operatorname{Re}P(\iu t) = t^4 +
1 \geq 1$। दायाँ अर्ध-समतल में शून्य:
अर्ध-चक्र $\{\abs z \leq
R,\ \operatorname{Re}z \geq 0\}$ की परिसीमा पर कोणांक सिद्धांत का उपयोग कीजिए। बड़े
चाप पर $\arg P
\approx \arg z^4$ $4\cdot\pi = 2\pi\cdot2$ भर घूमता है (चाप कोण $\pi$ फैलाता है)। $\iu R$ से
नीचे $-\iu R$ तक काल्पनिक अक्ष के अनुदिश: $P(\iu t) = (t^4 + 1) + 8\iu t$ दाएँ अर्ध-समतल में ही
रहता है ($\operatorname{Re} > 0$), अतः $\arg P$ $\intoo{-\pi/2}{\pi/2}$ के भीतर विचरता है और $R \to \infty$ पर शुद्ध
परिवर्तन $\to 0$ के साथ लौट आता है (दोनों अंत-बिंदु $\approx \arg t^4 = 0$)। कुल घूर्णन:
$\frac{4\pi + 0}
{2\pi} = 2$: अर्थात् दाएँ अर्ध-समतल में \emph{दो} शून्य --- जो संख्याओं
के संगत है: संयुग्म युग्म $\approx
1.04 \pm 1.73\iu$ का वास्तविक भाग धनात्मक है, और
वास्तविक शून्य $\approx -0.125$ तथा $\approx -1.96$ ऋणात्मक।
\end{solution}

\begin{solution}{pb:b3:residues:1}
\textbf{1.} $\sin(\pi z)$ के सरल शून्य ठीक $\Z$ पर हैं ($\sin\pi z = 0$ तभी और केवल
तभी जब $z \in \Z$, और वहाँ $(\sin\pi z)' = \pi\cos\pi
z \neq 0$), और $\cos(\pi n) \neq 0$: अतः $\pi\cot(\pi z) =
\pi\cos(\pi z)/\sin(\pi z)$ के $\Z$ पर सरल
ध्रुव हैं जिनमें
\[
\operatorname{Res}(\pi\cot\pi z,\ n)
= \frac{\pi\cos(\pi n)}{\pi\cos(\pi n)} = 1
\]
($g/h'$ नियम, \cref{met:b3:residues:compute})।

\textbf{2.} $\frac{\sin(\pi z)}{\pi z} = 1 - \frac{(\pi z)^2}6
+ \frac{(\pi z)^4}{120} - \cdots$ $0$ के निकट \omterm{def:b3:holomorphic:holo}{समविश्लेषिक} और शून्येतर है: अतः उसका
व्युत्क्रम \omterm{def:b3:holomorphic:holo}{समविश्लेषिक} है (\cref{def:b3:holomorphic:holo}: भागफल), जिसकी श्रेणी $1 +
\frac{(\pi z)^2}{6} + \frac{7(\pi z)^4}{360} + \cdots$ है
(पहचानिए: $u =
\frac{(\pi z)^2}6 - \frac{(\pi z)^4}{120}$ से $(1 - u)^{-1}$-शैली के गुणांक: $z^4$ गुणांक $\frac1{36} - \frac1{120} = \frac{7}{360}$ है)।
$\cos(\pi z) = 1 - \frac{(\pi z)^2}2 + \frac{(\pi
z)^4}{24} - \cdots$ से गुणा कीजिए और $z$ से भाग दीजिए:
\[
\pi\cot(\pi z) = \frac1z\Bigl[1 + \pi^2z^2\Bigl(\frac16 -
\frac12\Bigr) + \pi^4z^4\Bigl(\frac7{360} - \frac1{12} +
\frac1{24}\Bigr)\Bigr] + \cdots
= \frac1z - \frac{\pi^2}3\,z - \frac{\pi^4}{45}\,z^3 -
\cdots
\]
($\frac7{360} - \frac{30}{360} + \frac{15}{360} =
-\frac8{360} = -\frac1{45}$)।

\textbf{3.} ऊर्ध्वाधर भुजाएँ $z = \pm(N + \frac12) + \iu y$: $\cot$ की $\pi$-आवर्तिता से $\cot(\pi z) = \cot(\pm\frac\pi2
+ \iu\pi y) = -\tan(\iu\pi y) = -\iu\tanh(\pi y)$,
जिसका मापांक $\leq 1$ है। क्षैतिज भुजाएँ $z = x \pm \iu(N + \frac12)$: $\abs{\cot(a + \iu b)}^2 = \frac{\cos^2a +
\sinh^2b}{\sin^2a + \sinh^2b} \leq \frac{1 +
\sinh^2b}{\sinh^2b} = \coth^2 b$ से
\[
\abs{\cot(\pi z)} \leq \coth\bigl(\pi(N + \tfrac12)\bigr)
\leq \coth(\pi/2) < 1.1 .
\]
दोनों परिबंध $\leq 2$ हैं।

\textbf{4.} $C_N$ पर $F(z) =
\pi\cot(\pi z)f(z)$ पर \cref{thm:b3:residues:residue} लगाइए ($f$ के सभी ध्रुवों से
परे $N$):
\[
\frac1{2\iu\pi}\oint_{C_N}F = \sum_{n=-N}^{N}f(n) +
\sum_p\operatorname{Res}(F, p),
\]
जिसमें पूर्णांक ध्रुव $f(n)$ का योगदान देते हैं (प्रश्न 1; $f$ वहाँ
\omterm{def:b3:holomorphic:holo}{समविश्लेषिक})। $C_N$ पर: $\abs F \leq 2\pi\cdot
C\abs z^{-2} \leq 2\pi C N^{-2}$, और परिमाप $8(N +
\frac12)$ है: अतः समाकल $O(1/N) \to 0$
है। $N \to
\infty$ लीजिए: और दिखाया गया योग सूत्र मिल जाता है।

\textbf{5.} क्षय $\abs f = O(\abs z^{-2})$ ने \omterm{def:b3:holomorphic:contour}{परिरेखा समाकल} को मारा \emph{और} $\sum\abs{f(n)}$ को
अभिसरित कराया। $f(z) = \frac1{z + \frac12}$ के लिए: सममित योग $\sum_{-N}^N\frac1{n + \frac12}$ दूरबीनी होकर $0$ बन
जाते हैं (पद $n$ और $-n - 1$ निरस्त हो जाते हैं), और दायाँ पक्ष $-\operatorname{Res}\bigl(\frac{\pi\cot\pi z}{z + \frac12},
-\frac12\bigr) = -\pi\cot(-\frac\pi2) = 0$
है: संगत --- पर $\sum\abs{f(n)}$ अपसारित होता है; यह विधि केवल सममित (मुख्य
मान) सीमा परिकलित करती है।

\textbf{6.} $g(z) = \frac{\pi\cot(\pi z)}{z^2}$: शून्येतर पूर्णांकों पर \omterm{def:b3:residues:singularities}{अवशेष} $\frac1{n^2}$ वाले ध्रुव,
और $0$ पर, जहाँ प्रश्न 2 से
\[
g(z) = \frac1{z^3} - \frac{\pi^2}{3z} - \frac{\pi^4}{45}z -
\cdots :
\qquad \operatorname{Res}(g, 0) = -\frac{\pi^2}3 .
\]
$C_N$ पर \omterm{def:b3:holomorphic:contour}{परिरेखा समाकल} प्रश्न 4 की भाँति $0$ की ओर जाता है ($C_N$
पर $\abs{g} = O(N^{-2})$)। अतः $0 =
\sum_{n\neq0}\frac1{n^2} - \frac{\pi^2}3$: $2\zeta(2) =
\frac{\pi^2}3$, $\zeta(2) = \frac{\pi^2}6$।

\textbf{7.} $g(z) = \frac{\pi\cot(\pi z)}{z^4} = \frac1{z^5}
- \frac{\pi^2}{3z^3} - \frac{\pi^4}{45z} - \cdots$: $0$ पर \omterm{def:b3:residues:singularities}{अवशेष} $-\frac{\pi^4}{45}$ के बराबर, और $0 =
2\zeta(4) - \frac{\pi^4}{45}$: $\zeta(4) =
\frac{\pi^4}{90}$।

\textbf{8.} $g = \pi\cot(\pi z)/z^{2k}$ के साथ: $0$ पर \omterm{def:b3:residues:singularities}{अवशेष} $\pi\cot(\pi z)$ के प्रसार में $z^{2k-1}$
का गुणांक $a_{2k-1}$ है, और लुप्त होती परिरेखा $2\zeta(2k) + a_{2k-1} = 0$ दे देती है: $\zeta(2k) = -a_{2k-1}/2$,
जो $\pi^{2k}$ का कोई परिमेय गुणज है, क्योंकि कोटांजेंट के गुणांक परिमेय
हैं। एक और भाग-चरण $a_5 =
-\frac{2\pi^6}{945}$ देता है, जिससे $\zeta(6) = \frac{\pi^6}{945}$। $\zeta(3)$ के लिए: स्वाभाविक
अष्टि $g = \pi\cot(\pi z)/z^3$ विषमता से $\sum_{n\neq0}\frac1{n^3} = 0$ उत्पन्न कर देती है --- यह विधि $0 = 0$
सिद्ध करती है और विषम ज़ीटा मानों के प्रति संरचनात्मक रूप से अंधी
है ($\zeta(3)$ का कोई संवृत रूप ज्ञात नहीं; उसकी अपरिमेयता, अपेरी 1978,
के लिए बिलकुल भिन्न विचार चाहिए थे)।

\textbf{9.} $F(z) = \frac{\pi\cot(\pi z)}{(z - w)(z + w)}$ के पूर्णांकों पर ध्रुव हैं (\omterm{def:b3:residues:singularities}{अवशेष} $\frac1{n^2 - w^2}$, चिह्न पर
ध्यान दीजिए: $f(n) = \frac1{(n-w)(n+w)} = \frac1{n^2 -
w^2}$) और $\pm w$ पर प्रत्येक \omterm{def:b3:residues:singularities}{अवशेष} $\frac{\pi\cot(\pm\pi w)}{\pm2w} = \frac{\pi\cot(\pi
w)}{2w}$ वाले सरल ध्रुव
($\cot$ विषम है)। प्रश्न 4 का तर्क ($\abs f = O(\abs z^{-2})$)
\[
\sum_{n\in\Z}\frac1{n^2 - w^2} + \frac{\pi\cot(\pi w)}{w} =
0,
\qquad\text{अर्थात्}\qquad
\pi\cot(\pi w) = \frac1w + \sum_{n\geq1}\frac{2w}{w^2 -
n^2},
\]
देता है ($n = 0$ पद $-\frac1{w^2}$ है; $\pm n$ पुनःसमूहित कीजिए)। $\C\setminus\Z$ के संहतों
पर प्रसामान्य अभिसरण: $\abs w \leq
R$ और $n \geq 2R$ के लिए $\abs{\frac{2w}{w^2 - n^2}} \leq
\frac{2R}{n^2 - R^2} \leq \frac{8R}{3n^2}$।

\textbf{10.} $\abs w \leq r < 1$ के लिए: $\frac{2w}{w^2 - n^2} =
-\frac{2w}{n^2}\cdot\frac1{1 - w^2/n^2} =
-2\sum_{k\geq0}\frac{w^{2k+1}}{n^{2k+2}}$, जिसमें $\abs{\text{पद}} \leq 2r^{2k+1}/n^{2k+2}$, जो $(n, k)$ पर योगनीय
है: अतः श्रेणियों के लिए फुबिनी
\[
\pi\cot(\pi w) = \frac1w -
2\sum_{k\geq0}\zeta(2k+2)\,w^{2k+1} .
\]
का पुनर्विन्यास कर देती है। प्रश्न 2 से मिलान: $-2\zeta(2) = -\frac{\pi^2}3$ और $-2\zeta(4) = -\frac{\pi^4}{45}$ ---
वही मान। $\frac{\pi^2}6$ तक तीन मार्ग: \omterm{thm:b3:hilbert:parseval}{पार्सेवाल} (फूरिये श्रेणी), तार के ग्रीन
संकारक का संचिह्न, और कोटांजेंट के दो प्रसार; और यह कि आवृत्तियों
पर एक योग, एक संकारक संचिह्न और एक \omterm{def:b3:holomorphic:contour}{परिरेखा समाकल} सहमत हैं, कोई संयोग
नहीं --- प्रत्येक उसी स्पेक्ट्रमी सर्वसमिका का एक मुख है।

\textbf{11.} $R \geq 1$ स्थिर कीजिए और मान लीजिए $n_0$ सबसे छोटा पूर्णांक
$\geq 2R$ है। $\abs z \leq R$ और $n \geq n_0$ के लिए: $\abs{z^2/n^2} \leq \frac14$, अतः $1 - z^2/n^2 \in
\bar D(1, \frac14)$, जहाँ मुख्य लघुगणक
\omterm{def:b3:holomorphic:holo}{समविश्लेषिक} है, और
\[
\Bigl|\log\Bigl(1 - \frac{z^2}{n^2}\Bigr)\Bigr|
\leq \sum_{k\geq1}\frac1k\,\Bigl|\frac{z^2}{n^2}\Bigr|^k
\leq \frac{\abs{z^2/n^2}}{1 - \abs{z^2/n^2}}
\leq \frac{2R^2}{n^2} :
\]
जिसमें योग $S(z) = \sum_{n\geq n_0}\log(1 - z^2/n^2)$ $\bar D(0, R)$ पर प्रसामान्य रूप से अभिसरित होता है, जिसके
आंशिक योग $S_N$ \omterm{def:b3:holomorphic:holo}{समविश्लेषिक} हैं और एकसमान रूप से $\abs{S_N} \leq 2R^2\zeta(2)$। चूँकि $\abs{\eu^a - \eu^b} \leq
\eu^{\max(\abs a,\abs b)}\abs{a - b}$
(रेखाखंड पर माध्य मान परिबंध), अतः पुच्छ गुणनफल $\prod_{n_0\leq n\leq N} =
\eu^{S_N}$ $\bar D(0, R)$ पर एकसमान
रूप से शून्य-रहित $\eu^S$ पर अभिसरित होते हैं। स्थिर बहुपद $z\prod_{n<n_0}(1 - z^2/n^2)$ से
गुणा करने पर: $\bar D(0, R)$ पर एकसमान रूप से $P_N \to P$, और \cref{thm:b3:holomorphic:weierstrassconv} $P$ को वहाँ
\omterm{def:b3:holomorphic:holo}{समविश्लेषिक} बना देता है; और $R$ स्वेच्छ होने से $P$ समग्र है।
$\bar D(0, R)$ पर $P$ के शून्य बहुपद पूर्व-गुणक के शून्य हैं --- अर्थात्
मापांक $\leq R$ के पूर्णांक, प्रत्येक सरल ($(1 - z/n)(1 + z/n)$ के भिन्न सरल शून्य हैं,
$\eu^S$ का कोई नहीं): अतः $P$ का शून्य समुच्चय $\Z$ है, जिसमें सभी
शून्य सरल हैं।

\textbf{12.} परिमित गुणनफल का, उसके शून्यों से दूर, लघुगणकीय अवकलन:
\[
\frac{P_N'(z)}{P_N(z)} = \frac1z +
\sum_{n=1}^{N}\frac{-2z/n^2}{1 - z^2/n^2}
= \frac1z + \sum_{n=1}^{N}\frac{2z}{z^2 - n^2} .
\]
किसी \omterm{def:b3:topology:compact}{संहत} $K \subseteq \C\setminus\Z$ पर: एकसमान रूप से $P_N \to P$ और $P_N' \to P'$ (\cref{thm:b3:holomorphic:weierstrassconv}), तथा $\min_K
\abs P > 0$
($P$ केवल $\Z$ पर लुप्त होता है), अतः अंततः $\abs{P_N} \geq \frac12\min_K\abs P$ और $K$ पर
एकसमान रूप से $P_N'/P_N \to
P'/P$। बीच वाला सदस्य प्रश्न 9 से
$\frac1z + \sum_{n\geq1}\frac{2z}{z^2-n^2} = \pi\cot(\pi z)$
पर अभिसरित होता है: अतः $\C\setminus\Z$ पर $P'/P = \pi\cot(\pi z)$।

\textbf{13.} $\sin(\pi z)$ और $P$ एक ही शून्य समुच्चय $\Z$ वाले समग्र फलन
हैं, जिनके सभी शून्य सरल हैं (प्रश्न 1 और 11)। $m \in \Z$ के निकट $m$
पर \omterm{def:b3:holomorphic:holo}{समविश्लेषिक} और अलुप्त $\sigma, \psi$ के साथ $\sin(\pi z) = (z - m)\,\sigma(z)$ और $P(z) = (z - m)\,\psi(z)$ लिखिए (घात
श्रेणी का गुणनखंडन कीजिए): अतः $Q =
\sin(\pi z)/P = \sigma/\psi$ प्रत्येक पूर्णांक के पार
\omterm{def:b3:holomorphic:holo}{समविश्लेषिक} और शून्य-रहित रूप से बढ़ जाता है, और $\C\setminus\Z$ पर शून्य-रहित
फलनों के भागफल के रूप में शून्य-रहित है। वहाँ
\[
\frac{Q'}{Q} = \frac{(\sin\pi z)'}{\sin\pi z} -
\frac{P'}{P} = \pi\cot(\pi z) - \pi\cot(\pi z) = 0 ,
\]
अतः समग्र फलन $Q'$ $\C\setminus\Z$ पर लुप्त होता है, इसलिए \omterm{def:b3:topology:continuity}{सांतत्य} से सर्वत्र:
अतः $Q$ अचर है। $z \to
0$ पर: $\sin(\pi z)/z \to \pi$ और $P(z)/z \to 1$, अतः $Q =
\pi$:
\[
\sin(\pi z) = \pi z\prod_{n\geq1}
\Bigl(1 - \frac{z^2}{n^2}\Bigr) .
\]

\textbf{14.} $z = \frac12$ पर: $1 = \sin\frac\pi2 =
\frac\pi2\prod_{n\geq1}\bigl(1 - \frac1{4n^2}\bigr) =
\frac\pi2\prod\frac{4n^2-1}{4n^2}$, अतः
\[
\frac\pi2 = \prod_{n\geq1}\frac{4n^2}{4n^2 - 1}
= \lim_{N\to\infty}\prod_{n=1}^N
\frac{(2n)(2n)}{(2n-1)(2n+1)}
= \lim_{N\to\infty}
\frac{2\cdot2\cdot4\cdot4\cdots(2N)(2N)}
{1\cdot3\cdot3\cdot5\cdots(2N-1)(2N+1)} :
\]
अर्थात् वालिस का गुणनफल, ऑयलर के गुणनखंडन का एक-पंक्ति उपप्रमेय।

\textbf{15.} $\abs z \leq r < 1$ के लिए प्रत्येक गुणक $D(1, r^2) \subseteq D(1, 1)$ में है, अतः $P(z)/z = \exp\bigl(
\sum_{n\geq1}\log(1 - z^2/n^2)\bigr)$:
प्रत्येक आंशिक गुणनफल किसी आंशिक योग का चरघातांक है, और दोनों पक्ष
$\exp$ की \omterm{def:b3:topology:continuity}{सांतत्य} से सीमा तक चले जाते हैं। द्विक श्रेणी
\[
\sum_{n\geq1}\log\Bigl(1 - \frac{z^2}{n^2}\Bigr)
= -\sum_{n\geq1}\sum_{k\geq1}\frac{z^{2k}}{k\,n^{2k}}
= -\sum_{k\geq1}\frac{\zeta(2k)}k\,z^{2k}
= -\zeta(2)\,z^2 + O(z^4)
\]
श्रेणियों के लिए फुबिनी से पुनर्विन्यस्त हो जाती है: $\sum_{n,k}
\frac{r^{2k}}{kn^{2k}} \leq \sum_k\zeta(2k)r^{2k} \leq
\zeta(2)\frac{r^2}{1-r^2} < \infty$। अतः
\[
P(z) = z\,\exp\bigl(-\zeta(2)z^2 + O(z^4)\bigr)
= z - \zeta(2)\,z^3 + O(z^5) ,
\]
और प्रश्न 13 इसकी तुलना $\sin(\pi z) = \pi z -
\frac{\pi^3}6z^3 + O(z^5)$ से करता है: $\pi\zeta(2) = \frac{\pi^3}6$, अर्थात् $\zeta(2) = \frac{\pi^2}6$। योगात्मक
मुख (प्रश्न 10) और गुणनात्मक मुख एक ही संख्या परिकलित करते हैं।

\textbf{16.} किसी \omterm{def:b3:topology:compact}{संहत} $K \subseteq \C\setminus\Z$ पर आंशिक योग $S_N = \frac1z + \sum_{n\leq N}\bigl(
\frac1{z-n} + \frac1{z+n}\bigr)$ (प्रश्न 9, पदों को
$\frac{2z}{z^2-n^2} = \frac1{z-n} +
\frac1{z+n}$ के रूप में पुनःसमूहित करके) एकसमान रूप से $\pi\cot(\pi z)$ पर अभिसरित होते
हैं, अतः \cref{thm:b3:holomorphic:weierstrassconv} $K$ पर एकसमान रूप से $S_N' \to
(\pi\cot\pi z)' = -\pi^2/\sin^2(\pi z)$ दे देता है। चूँकि $S_N' = -\sum_{\abs n\leq N}(z - n)^{-2}$:
\[
\frac{\pi^2}{\sin^2(\pi z)} =
\sum_{n\in\Z}\frac1{(z - n)^2} ,
\]
जिसका अभिसरण $\C\setminus\Z$ के संहतों पर प्रसामान्य है (पद $O(n^{-2})$)।

\textbf{17.} $z = \frac12$ पर बायाँ पक्ष $\pi^2$ है; और दाईं ओर $(\frac12 - n)^2 = \frac{(2n-1)^2}4$, जिसमें
जैसे-जैसे $n$ $\Z$ पर चलता है $2n -
1$ ठीक एक बार सभी विषम पूर्णांकों
पर चलता है:
\[
\pi^2 = \sum_{n\in\Z}\frac{4}{(2n-1)^2}
= 8\sum_{m\geq0}\frac1{(2m+1)^2},
\qquad
\sum_{m\geq0}\frac1{(2m+1)^2} = \frac{\pi^2}8 .
\]
सम-विषमता विभाजन: $\zeta(2) = \frac{\pi^2}8 +
\sum_{n\geq1}\frac1{(2n)^2} = \frac{\pi^2}8 +
\frac{\zeta(2)}4$, अतः $\frac34\zeta(2) = \frac{\pi^2}8$ और एक बार फिर $\zeta(2) = \frac{\pi^2}6$।

\textbf{18.} $c = \cot\theta$ और $\cot(2\theta) =
\frac{c^2-1}{2c}$ के साथ: $\cot\theta - 2\cot(2\theta) = c -
\frac{c^2-1}c = \frac1c = \tan\theta$। अतः $\pi\tan(\pi z)
= \pi\cot(\pi z) - 2\pi\cot(2\pi z)$, और $z$ पर
तथा $2z$ पर प्रश्न 9
\[
\pi\cot(\pi z) = \frac1z +
\sum_{n\geq1}\frac{2z}{z^2 - n^2},
\qquad
2\pi\cot(2\pi z) = \frac1z +
\sum_{n\geq1}\frac{8z}{4z^2 - n^2} .
\]
दे देता है। दोनों श्रेणियाँ ध्रुवों से बाहर प्रत्येक स्थिर $z$ पर
निरपेक्ष रूप से अभिसरित होती हैं, अतः अंतर को इच्छानुसार पुनःसमूहित
किया जा सकता है: दूसरी श्रेणी में सम पद $n = 2m$ $\frac{8z}{4z^2 -
4m^2} = \frac{2z}{z^2 - m^2}$ देते हैं और पहली
श्रेणी को पूरी तरह निरस्त कर देते हैं, जिससे
\[
\pi\tan(\pi z) = -\sum_{m\geq0}\frac{8z}{4z^2 - (2m+1)^2}
= \sum_{m\geq0}\frac{8z}{(2m+1)^2 - 4z^2} ,
\]
बचता है, जो $\frac12 + \Z$ से बचने वाले संहतों पर प्रसामान्य है (पद $O(m^{-2})$)।

\textbf{19.} $\abs z \leq r < \frac12$ के लिए:
\[
\frac{8z}{(2m+1)^2 - 4z^2} = \frac{8z}{(2m+1)^2}
\sum_{k\geq0}\Bigl(\frac{4z^2}{(2m+1)^2}\Bigr)^{k},
\]
जिसमें $\sum_{m,k}8r\,(4r^2)^k(2m+1)^{-2k-2} < \infty$, क्योंकि $4r^2 < 1$: अतः फुबिनी द्विक योग का पुनर्विन्यास
\[
\pi\tan(\pi z) =
\sum_{k\geq0}8\cdot4^k\,\lambda(2k+2)\,z^{2k+1} .
\]
में कर देती है। $\pi\tan(\pi z) = \pi^2z + \frac{\pi^4}3z^3 +
O(z^5)$ के सापेक्ष: $z$ का गुणांक $8\lambda(2) = \pi^2$ देता है ---
फिर से प्रश्न 17 --- और $z^3$ का गुणांक $32\lambda(4) = \frac{\pi^4}3$, अर्थात् $\lambda(4) =
\frac{\pi^4}{96}$। सम हर
हटाने पर $\lambda(4) = \zeta(4) - 2^{-4}\zeta(4) =
\frac{15}{16}\zeta(4)$: $\zeta(4) = \frac{16}{15}\cdot
\frac{\pi^4}{96} = \frac{\pi^4}{90}$, जो प्रश्न 7 से मेल खाता है।

\textbf{20.} $\cot\frac\theta2 - \cot\theta =
\frac{\cos\frac\theta2\,\sin\theta -
\cos\theta\,\sin\frac\theta2}{\sin\frac\theta2\,\sin\theta}
= \frac{\sin\frac\theta2}{\sin\frac\theta2\,\sin\theta} =
\frac1{\sin\theta}$, जिसका अंश $\sin(\theta -
\frac\theta2)$ है। $\theta = \pi z$ के साथ $\frac z2$ पर प्रश्न
9 $\pi\cot\frac{\pi z}2 = \frac2z +
\sum_{n\geq1}\frac{4z}{z^2 - 4n^2}$ के रूप में पढ़ा जाता है, अतः
\[
\frac{\pi}{\sin(\pi z)}
= \pi\cot\frac{\pi z}2 - \pi\cot(\pi z)
= \frac1z + \sum_{n\geq1}\frac{4z}{z^2 - 4n^2}
- \sum_{n\geq1}\frac{2z}{z^2 - n^2} .
\]
निरपेक्ष अभिसरण सम-विषमता पुनःसमूहन की अनुमति देता है: घटाई गई श्रेणी
में सम $n
= 2m$ $\frac{2z}{z^2-4m^2}$ का योगदान देता है, जिससे पहले योग में से $+\frac{2z}{z^2-4m^2}$
बचता है, जबकि विषम $n$ चिह्न $-$ के साथ बच जाता है:
\[
\frac{\pi}{\sin(\pi z)} = \frac1z +
\sum_{n\geq1}(-1)^n\,\frac{2z}{z^2 - n^2} .
\]
चिह्न जाँच: $\operatorname{Res}\bigl(\pi/\sin(\pi z),
n\bigr) = \pi/(\pi\cos\pi n) = (-1)^n$ (\cref{met:b3:residues:compute}), और $(-1)^n\frac{2z}{z^2 -
n^2} = \frac{(-1)^n}{z-n} + \frac{(-1)^n}{z+n}$ $\pm n$ पर ठीक वही \omterm{def:b3:residues:singularities}{अवशेष} ढोता है।

\textbf{21.} $\abs z \leq r < 1$ के लिए प्रत्येक पद का प्रश्न 19 की भाँति प्रसार
करके ($(-1)^n\frac{2z}{z^2-n^2} =
(-1)^{n-1}\frac{2z}{n^2}\sum_k\frac{z^{2k}}{n^{2k}}$) और फुबिनी लगाकर:
\[
\frac{\pi}{\sin(\pi z)} = \frac1z +
2\sum_{k\geq0}\eta(2k+2)\,z^{2k+1},
\qquad
\eta(s) = \sum_{n\geq1}\frac{(-1)^{n-1}}{n^s} .
\]
टेलर की ओर: $\sin(\pi z) = \pi z(1 - \frac{(\pi z)^2}6 +
O(z^4))$ $\frac\pi{\sin\pi z} = \frac1z +
\frac{\pi^2}6z + O(z^3)$ देता है: $2\eta(2) = \frac{\pi^2}6$, अतः $\eta(2) = \frac{\pi^2}{12}$। संगति: $\eta(2) =
\zeta(2) - 2\sum_n(2n)^{-2} = (1 - 2^{1-2})\zeta(2) =
\frac{\zeta(2)}2 = \frac{\pi^2}{12}$।

\textbf{22.} $z = \frac12$ पर प्रश्न 20
\[
\pi = 2 + \sum_{n\geq1}(-1)^n\frac{1}{\frac14 - n^2}
= 2 + 4\sum_{n\geq1}\frac{(-1)^{n-1}}{(2n-1)(2n+1)} .
\]
देता है। $\frac1{(2n-1)(2n+1)} = \frac12\bigl(\frac1{2n-1} -
\frac1{2n+1}\bigr)$ के साथ, और दोनों एकांतर श्रेणियों के अभिसारी होने से
(लाइब्नित्स कसौटी), योग $\frac12\bigl[L -
(1 - L)\bigr] = L - \frac12$ के रूप में बँट जाता है, जहाँ $L = 1 - \frac13 +
\frac15 - \cdots$ और
$\frac13 - \frac15 + \frac17 - \cdots
= 1 - L$। अतः $\pi = 2 + 4(L - \frac12) = 4L$:
\[
\frac\pi4 = 1 - \frac13 + \frac15 - \frac17 + \cdots
\]
शिक्षा: प्रत्येक अष्टि \omterm{def:b3:residues:singularities}{मेरोमॉर्फिक} है, जिसके ध्रुव किसी समांतर श्रेढ़ी
पर हैं और \omterm{def:b3:residues:singularities}{अवशेष} निर्धारित हैं। कोटांजेंट प्रत्येक पूर्णांक पर \omterm{def:b3:residues:singularities}{अवशेष}
$1$ रखता है और $f(n)$ का योग लेता है; उसका अवकलज ध्रुवों का वर्ग
कर देता है और $\lambda(2)$ का मूल्य आँकता है; टैंजेंट ध्रुवों को $\frac12 + \Z$ पर
हटा देता है और विषम हरों का मूल्य आँकता है; और $\pi/\sin$ पूर्णांक ध्रुव
रखे रहता है पर \omterm{def:b3:residues:singularities}{अवशेष} $(-1)^n$ को एकांतरित कर देता है, जिससे एकांतर
श्रेणी मिलती है। प्रथम-कोटि की सभी अष्टियाँ विषम फलन हैं: $n$ को
$-n$ के साथ युग्मित करने से सम-घात गुणांक दुगुने हो जाते हैं और
विषम गुणांक नष्ट, अतः $\zeta(2k)$ यंत्रवत बह निकलता है जबकि $\zeta(3)$ कभी
प्रकट नहीं होता। यह अंधापन सम-विषमता का है, तकनीक की कमी का नहीं।

\textbf{23.} $w = 2\iu\pi z$ के साथ:
\[
\iu\pi z + \frac{2\iu\pi z}{\eu^{2\iu\pi z} - 1}
= \iu\pi z\,\frac{\eu^{2\iu\pi z} + 1}{\eu^{2\iu\pi z} - 1}
= \iu\pi z\,\frac{\eu^{\iu\pi z} +
\eu^{-\iu\pi z}}{\eu^{\iu\pi z} - \eu^{-\iu\pi z}}
= \pi z\,\frac{\cos\pi z}{\sin\pi z} = \pi z\cot(\pi z) .
\]
अतः $w = 2\iu\pi z$ पर जनक फलन का उपयोग करते हुए (और $\iu\pi z$ पद को निरस्त करते
$B_1 = -\frac12$ के साथ, जिससे आगे विषम $B$ लुप्त हो जाते हैं):
\[
\pi z\cot(\pi z) = \sum_{k\geq0}\frac{B_{2k}}{(2k)!}
(2\iu\pi z)^{2k}
= 1 + \sum_{k\geq1}(-1)^k\frac{(2\pi)^{2k}B_{2k}}{(2k)!}
z^{2k} .
\]
दूसरी ओर आंशिक भिन्न प्रसार (भाग IV) $\pi z\cot(\pi z) = 1 - 2\sum_{k\geq1}\zeta(2k)z^{2k}$ देता है (प्रत्येक $\frac{2z^2}{z^2 - n^2} =
-2\sum_k\frac{z^{2k}}{n^{2k}}$
का प्रसार कीजिए और $n$ पर योग लीजिए; $\abs z < 1$ के लिए प्रसामान्य अभिसरण
अदला-बदली को उचित ठहराता है)। गुणांकों की तुलना: $-2\zeta(2k) =
(-1)^k\frac{(2\pi)^{2k}B_{2k}}{(2k)!}$, जो कहा गया
सूत्र है।

\textbf{24.} $k = 1$: $\frac{(2\pi)^2}{2\cdot2}\cdot\frac16
= \frac{\pi^2}6$। $k = 2$: $-\frac{(2\pi)^4}{2\cdot24}
\cdot\bigl(-\frac1{30}\bigr) = \frac{16\pi^4}{48\cdot30} =
\frac{\pi^4}{90}$। $k = 3$: $\frac{(2\pi)^6}{2\cdot720}
\cdot\frac1{42} = \frac{64\pi^6}{1440\cdot42} =
\frac{\pi^6}{945}$।

\textbf{25.} $C(z) = \pi z\cot(\pi z) = 1 -
2\sum_{k\geq1}\zeta(2k)z^{2k}$ लिखिए (प्रश्न 23)। $(\cot u)'
= -1 - \cot^2u$ का उपयोग करते हुए $C = \pi z\cot(\pi z)$
का सीधा अवकलन:
\[
zC'(z) = \pi z\cot(\pi z) - \pi^2z^2\bigl(1 +
\cot^2(\pi z)\bigr) = C - \pi^2z^2 - C^2 .
\]
अब दोनों पक्षों का $z^2$ की घातों में प्रसार कीजिए। बायाँ पक्ष:
$\sum_k(-4k)\,\zeta(2k)\,z^{2k}$। दायाँ पक्ष: श्रेणी का वर्ग करने पर
\[
C - C^2 = 2\sum_{k\geq1}\zeta(2k)z^{2k}
- 4\sum_{k\geq2}\Bigl(\sum_{j=1}^{k-1}\zeta(2j)
\zeta(2k-2j)\Bigr)z^{2k},
\]
और पद $-\pi^2z^2 = -6\zeta(2)z^2$ केवल $k =
1$ को समायोजित करता है। $k \geq 2$ के लिए $z^{2k}$ के
गुणांकों की तुलना करने पर:
\[
-4k\,\zeta(2k) = 2\,\zeta(2k) -
4\sum_{j=1}^{k-1}\zeta(2j)\,\zeta(2k-2j),
\]
अर्थात् $\bigl(k + \frac12\bigr)\zeta(2k) =
\sum_{j=1}^{k-1}\zeta(2j)\zeta(2k-2j)$। ($k = 1$ पर यह सर्वसमिका $-4\zeta(2) = 2\zeta(2) - 6\zeta(2)$ के रूप में पढ़ी जाती है:
एक संगति जाँच, नई सूचना नहीं।) अनुप्रयोग: $k =
2$: $\frac52\zeta(4) = \zeta(2)^2 = \frac{\pi^4}{36}$, अतः $\zeta(4) = \frac{\pi^4}{90}$;
$k = 3$: $\frac72\zeta(6) =
2\zeta(2)\zeta(4) = \frac{\pi^6}{270}$, अतः $\zeta(6) =
\frac{2}{7}\cdot\frac{\pi^6}{270} = \frac{\pi^6}{945}$। एक अतिक्रांत बीज ($\zeta(2) = \frac{\pi^2}6$), और विशुद्ध बीजगणित
प्रत्येक सम ज़ीटा मान उत्पन्न कर देता है।
\end{solution}
